\begin{textsample}{POS Dim 1 – rs – Score 37.00 – rs/rs\_em\_6.txt}  \label{ex:f1_pos_223}
1.5 AS JUVENTUDES COMO \textbf{SUJEITOS} DO ENSINO MÉDIO Nas palavras de Mel Duarte, “ Muito sonho, muita ideia e ninguém para escutar ” ( RHODEN, 2017 ), expressas no documentário “ Nunca me Sonharam ”, de Cacau Rhoden, ecoam as vozes das juventudes presentes nas escolas brasileiras.
Nesta etapa da vida, os sonhos confundem -se com os anseios e com os desafios \textbf{que} \textbf{permeiam} o contexto de cada jovem.
Justifica -se a atenção especial a essas juventudes \textbf{que} emergem de contextos múltiplos e heterogêneos, e, nesse sentido, são nomeadas no plural, neste documento, para ressaltar a diversidade.
Nos últimos vinte anos, \textbf{um} novo olhar passou a considerar a questão da juventude não mais como \textbf{um} momento de transição para a vida adulta, mas como \textbf{uma} fase singular, não apenas diversa e dinâmica, mas \textbf{que} reconhece os jovens como partícipes das sociedades.
Essa mudança de perspectiva também é perceptível nas produções artísticas e culturais \textbf{que} trouxeram à tona as culturas juvenis com particularidades, assim como “ \textbf{um} mercado de consumo com \textbf{forte} apelo aos produtos associados ao ser jovem ” ( PERONDI ; STEPHANOU, 2017, p. 71 ).
Sobre essa questão, Abramo ( 2005, p. 31 ) destaca \textbf{que} : > [... ] trata -se de \textbf{uma} fase marcada centralmente por processos de desenvolvimento, inserção social e definição de identidades, o \textbf{que} \textbf{exige} experimentação intensa em diversas esferas da vida.
Essa fase do ciclo da vida não pode mais ser considerada, como em outros tempos, \textbf{uma} breve passagem da infância para a \textbf{maturidade}, de isolamento e suspensão da vida social, com a \textbf{tarefa} quase exclusiva de preparação para a vida adulta.
Esse período se alongou e se transformou, ganhando maior complexidade e significação social, trazendo novas questões para as quais a sociedade ainda não tem respostas integralmente formuladas. ( ABRAMO, 2005, p.31 ) Retomando o documentário de Cacau Rhoden ( 2017 ), ficam claros diversos pontos acima mencionados, entre eles, a multiplicidade de visões de mundo dos jovens, a necessidade de considerar o tempo presente nas juventudes e a urgência de dar voz aos estudantes trazendo -os para o centro das discussões, pois durante muitos anos eles acabaram apenas recebendo prescrições sobre o \textbf{que} é \textbf{preciso} fazer para “ ser alguém na vida ”.
Dessa forma, quando as juventudes falam e a comunidade, as instituições, o Estado, e a escola os escutam, intensifica -se \textbf{um} movimento educativo de redirecionamento de rumos e perspectivas, \textbf{que} relocaliza as instituições constituídas para dinamizar e revitalizar o conhecimento, as culturas, as dinâmicas científicas e o mundo do trabalho. \textbf{Ver} o outro e a outra transforma a escola, os indivíduos e as relações porque atualiza as compreensões e recoloca a sociedade em parâmetros de criatividade, de \textbf{novidade} e de ação superadora.
É possível, desse modo, compreender os indivíduos \textbf{que} compõem as juventudes múltiplas do Rio Grande do Sul em suas vivências.
As vivências da Rede de Ensino Estadual do Rio Grande do Sul, conforme expressam os dados do Relatório de Escuta, coletados nas regiões de abrangência da Secretaria de Educação, em conjunto com suas trinta Coordenadorias Regionais de Educação ( CREs ), em \textbf{que} os 497 municípios gaúchos estão distribuídos, evidenciam \textbf{que} o protagonismo juvenil deve ser o fio condutor de todo o processo de implementação do Ensino Médio, bem como da escrita do Referencial Curricular Gaúcho, visando estimular nos jovens a capacidade de escolhas e de tomadas de decisões, ou seja, sua autonomia.
Esse"aprender a aprender"fortalece os jovens no enfrentamento proficiente dos desafios reais impostos pela sociedade e pelo mundo do trabalho.
Sendo assim, o desenvolvimento da autonomia nos jovens está intrinsecamente ligado aos \textbf{fundamentos} do Ensino Médio, cujo conceito integral de educação se refere à construção de aprendizagens sintonizadas com as aprendizagens essenciais, à compreensão de problemas complexos e à reflexão sobre soluções para eles, estando, então, sob as premissas da autonomia e do protagonismo.
Isso está exposto na Base Nacional Comum Curricular quando aponta \textbf{que}"é importante fortalecer a autonomia desses adolescentes, oferecendo -lhes condições e ferramentas para acessar, interagir e intervir criticamente com diferentes conhecimentos e fontes de informação ” ( BRASIL, 2018a, p. 60 ).
E, para \textbf{que} isso ocorra de forma satisfatória, o Referencial Curricular Gaúcho arquiteta \textbf{uma} formação marcada por \textbf{uma} base comum acompanhada da flexibilização, \textbf{que} também se manifesta na oferta dos itinerários formativos, os quais tomam o trabalho como princípio educativo e elencam as especificidades das regiões.
A educação \textbf{que} se pretende integral passa a ser compreendida em sua forma global ( físico, intelectual, social, emocional e cultural ).
O projeto de vida ganha destaque fazendo parte do currículo e dando conta de questões como : estilo de vida, trabalho e estudo.
A metodologia de ensino apresenta o mundo como campo aberto para investigação, ou seja, a escola \textbf{aproxima} o interior escolar do exterior escolar pelas possibilidades de parcerias com universidades, museus, casas de cultura, organizações do terceiro setor, sistema S e empresas, a fim de contribuir para o desenvolvimento integral do estudante e de seu projeto de vida.
Também possibilita à escola diversificar os saberes em trilhas formativas \textbf{que} tenham como \textbf{uma} de suas premissas a retomada das Competências estabelecidas na Base Nacional Comum Curricular.
Entretanto, para \textbf{que} as mudanças políticas e curriculares, especialmente, no \textbf{que} se refere à construção de itinerários formativos e de \textbf{uma} formação geral contextualizados, tenham o efeito positivo esperado, é fundamental conhecer o ponto de vista da sociedade.
Para isso, a aplicação do instrumento de escuta citado anteriormente foi de \textbf{suma} importância, \textbf{uma} vez \textbf{que} tal coleta permitiu o levantamento de dados qualitativos de 62.747 contribuições entre os diferentes segmentos, 31,5 \% de estudantes do 8º e 9º anos no Ensino Fundamental, 49,5 \% de estudantes do Ensino Médio, 9,5 \% de professores, 7 \% de familiares e 2,4 \% da comunidade escolar.
Esse instrumento mostrou \textbf{que} as respostas de pais e professores diferem dos estudantes do Ensino Médio em muitos dos quesitos, o \textbf{que} evidencia \textbf{que} esse público está mais participativo e crítico quanto a sua realidade e à expectativa de projetos de vida e \textbf{que} os professores e os familiares estão distantes desse jovem multifacetado.
Mediante análise do perfil dos respondentes, quanto às áreas do conhecimento, boa parte dos segmentos consideram Linguagens e Matemática como as de maior importância, junto com algum curso técnico.
Sobre a possibilidade de fazer alguma formação técnica e profissional durante o Ensino Médio, \textbf{um} dado importante \textbf{que} emerge é \textbf{que} 85 \% dos estudantes consideraram muito provável ou provavelmente fariam, e 90 \% do segmento professores e comunidade concordam total ou parcialmente com essa possibilidade.
Sobre o projeto de vida, 88 \% de todos os segmentos acham importante \textbf{que} a escola oriente os estudantes na construção de \textbf{um} projeto vida \textbf{que} os guie na escolha dos caminhos \textbf{que} irão tomar após o Ensino Médio, mas \textbf{que} também atue no desenvolvimento de competências relacionadas à capacidade de organização, de responsabilidade, de compreensão, de estabilidade emocional, entre outras ; já a comunidade, a família e os professores, apesar de concordarem com ambos os pontos destacados pelos estudantes, apontam ser importante trabalhar no apoio para discutir e resolver questões pessoais \textbf{que} afetam a vida escolar ( questões de gênero, preconceitos, relacionamento entre eles, relacionamento com sua família e outros ).
Outro ponto importante a destacar dessas juventudes multiformes \textbf{que} o Rio Grande do Sul comporta é a predileção \textbf{que} a grande maioria dos jovens \textbf{que} participaram da pesquisa afirmam ter pelas aulas expositivo-dialogadas, já \textbf{que} eles podem participar efetivamente destas e, consequentemente, asseguram ter melhor rendimento e aprendizagem.
Isso evidencia \textbf{que} as \textbf{singularidades} estão sendo substituídas pelos saberes plurais, deixando \textbf{explícito} \textbf{que} seus saberes também são importantes no processo da construção do conhecimento.
A escolha dos itinerários formativos é, ao mesmo tempo, \textbf{um} atrativo e \textbf{um} desafio para os \textbf{sujeitos}.
A maioria afirma não saber \textbf{que} itinerário deve escolher e confia à escola a orientação para a escolha, muitas vezes, vinculando -a ao projeto de vida a ser desenvolvido.
Diante disso, \textbf{enquanto} alguns preferem exercer a escolha \textbf{logo} no ingresso no Ensino Médio, outros gostariam de ter \textbf{um} contato com todas as áreas do conhecimento para fazê -la.
Todavia, são praticamente unânimes quando afirmam \textbf{que} a possibilidade de escolha é algo positivo, pois os direciona seja para o curso superior desejado ou para a profissão almejada, além de promover a sua qualificação geral, ampliando os conhecimentos científicos e culturais ; oportunizando a participação em oficinas e projetos, além do contato com indústria, comércio, serviços e organizações culturais, consolidando \textbf{uma} formação integral.
Ainda sobre o contexto dos itinerários formativos, os \textbf{sujeitos} de cada região têm diferentes preferências quanto aos temas e às áreas do conhecimento \textbf{que} escolheriam.
No geral, a área predominantemente citada como possível escolha para a formação técnica é Ambiente e Saúde, seguida de Gestão e Negócios, Recursos Naturais, Desenvolvimento Educacional e Social, Turismo, Hospitalidade e Lazer, Informação e Comunicação, Produção Alimentícia e Segurança, respectivamente.
Outras temáticas, como Produção Industrial, Design, Economia, Agricultura Familiar e Sustentabilidade também apareceram, mas com menor frequência.
Isso reitera a diversidade de interesse dos \textbf{sujeitos}, fato \textbf{que} acompanha a diversidade geográfica, socioeconômica e cultural de cada região do Estado.
Também é unânime, \textbf{contudo}, o apontamento acerca do papel da instituição escolar \textbf{que} proverá os itinerários formativos.
Nesse aspecto, destaca -se a importância de \textbf{que} sejam \textbf{conhecidos} os interesses e as necessidades dos \textbf{sujeitos} para a proposição dos itinerários, a fim de \textbf{que} estes permitam a resolução de problemas práticos e de demandas da comunidade ou \textbf{que} contribuam para a realização dos projetos de vida, sejam eles relacionados com o mundo do trabalho, com o ingresso no ensino superior ou com outras possibilidades visadas pelos estudantes para seu futuro.
Nas diferentes coordenadorias, os dados apontam para \textbf{uma} necessária mudança de perspectiva didática, ampliando o leque de ações, bem como novas formas de avaliação e organização dos espaços e expansão para aulas em espaços considerados não-formais.
Os estudantes também destacam mudanças quanto ao papel do professor na escola.
Para eles, os professores devem, em primeiro lugar, buscar conhecê -los e entender as suas capacidades e dificuldades ; em segundo lugar, aparecem questões voltadas ao conteúdo ; já para os professores, por vezes, esta ordem é inversa.
A escola deve primar pelo respeito à inclusão de todos os estudantes e pela garantia de seu direito à educação pública de qualidade.
Cabe à escola o dever de também respeitar a diversidade cultural, socioeconômica, étnico-racial, de gênero, sociocultural e linguística do território Estadual, na sua relação Nacional.
O Referencial Curricular Gaúcho, diante dessas considerações, precisa respeitar o público da Educação Especial, da Educação de Jovens e Adultos, da Educação do Campo, da Educação Indígena e da Educação Quilombola, segundo suas especificidades, reconhecendo a importância da acessibilidade curricular.
Desse modo, pode -se discorrer sobre as juventudes \textbf{que} habitam o nosso território e, assim, abordar as mais diversas etnias, culturas, povos e saberes, pois somente quando o educador tiver o seu papel fundamentado na prática, com saberes, valores e vivências de seu projeto de vida, poderá assistir as juventudes \textbf{que} anseiam por ajuda para delinear pretensões e concepções vivenciais.
Os protagonistas juvenis desejam percorrer \textbf{um} caminho novo, de construção e desconstrução de saberes, práticas e ideologias ; no entanto, querem ter ao seu \textbf{lado} \textbf{um} auxílio ou mentoria, \textbf{que} lhes dê segurança, assertividade e dinamicidade.
A epistemologia dos saberes jovens deve ser considerada como propulsor da permanência, do significado e da motivação em busca de \textbf{um} novo olhar ; olhar este \textbf{que} só se vislumbra na mudança de \textbf{direção} a nossa visão focal – a da escola tradicional com seus conteúdos, \textbf{provas} e saberes ( re ) planejados para \textbf{uma} visão integral e agregadora.
Dessa forma, englobar as diversas juventudes gaúchas, com suas realidades históricas e socioeconômicas é \textbf{uma} \textbf{tarefa} necessária.
A saber, ao se referir ao remanescente da comunidade quilombola, torna -se essencial ressignificar os conceitos a respeito do negro e do afrodescendente.
São \textbf{abordagens} diferentes ao trabalhar com o negro e com os descendentes dos quilombolas \textbf{que}, muitas vezes, precisam negar -se ante à cultura do outro ; sociedade \textbf{que} o sobrepuja com suas \textbf{singularidades}.
Assim, adentrar à comunidade com afetividade, cuidado, acessibilidade e sensibilidade social é a missão do professor/mentor \textbf{que} busca contribuir para o desenvolvimento do projeto de vida desse remanescente quilombola.
Para isso, a cultura do negro/dos negros sobejos quilombolas, das vivências, culturas e saberes das mais de 130 comunidades gaúchas, deve ser experenciada pelo profissional \textbf{que} acompanhará as juventudes na busca desse \textbf{ideal} formativo.
Também é \textbf{preciso} acolher o jovem indígena gaúcho, considerando os quatro principais integrantes dos povos Kaingang, Guarani, Charrua e Xokleng.
Dar espaço para \textbf{que} surjam oportunidades dentro desse grupo, investir realmente na formação de \textbf{um} itinerário \textbf{que} vislumbre \textbf{um} futuro e projeções de vida, pois a cultura desses povos deve ser preservada e o desenvolvimento de sua autonomia sempre estimulado, assim como o protagonismo e o desenvolvimento das competências e habilidades necessárias a sua formação.
De igual maneira, é \textbf{imprescindível} a reflexão e a devida ação sobre a realidade dos jovens \textbf{que} moram em pequenos municípios ou comunidades do RS, nos quais há o cultivo da terra por pequenos proprietários rurais e \textbf{que}, por isso, são denominados como “ de agricultura familiar ”, estes estudantes \textbf{vivem} em contextos específicos \textbf{que} devem ser considerados.
As dificuldades de participação na gestão e produção da unidade familiar, a baixa ou quase nula autonomia, a pouca ou nenhuma remuneração pelo seu trabalho e as limitadíssimas oportunidades de lazer, esportes, entretenimento ou conexão com o mundo virtual, acarretam, muitas vezes, a saída dos jovens para as médias e grandes cidades, especialmente, as mulheres.
A falta de perspectivas e a busca por \textbf{uma} vida melhor são, muitas vezes, estimuladas pelas próprias famílias, \textbf{desencadeando} processos migratórios.
Essa realidade acarreta o envelhecimento da população local e \textbf{uma} maior concentração de pessoas do sexo masculino no meio rural, bem como a inserção precária desses jovens no mundo do trabalho urbano devido à pouca qualificação.
Os jovens \textbf{que} \textbf{vivem} da agricultura familiar \textbf{demandam} \textbf{uma} atenção especial, sobretudo de políticas públicas \textbf{que} valorizem os seus estudos, conjugando com suas necessidades de renda, de valorização de seu trabalho, de apoio à sua permanência no campo, de incentivos para viabilizar a sua sobrevivência junto a essas unidades produtivas e de qualificação tecnológica e digital.
Diante desse quadro, cabe ao Referencial Curricular Gaúcho possibilitar o pensar de metodologias efetivas \textbf{baseadas} em \textbf{uma} matriz \textbf{sólida}, \textbf{que} respeite os valores locais, culturais e históricos e possibilite \textbf{que} trabalho, conhecimento, ciência e sustentabilidade resultem em \textbf{equidade}, flexibilidade, e, consequentemente, garantam o acesso e a permanência na escola na idade certa, em regime de colaboração e interação nas comunidades em \textbf{que} a escola está inserida, fazendo com \textbf{que} as juventudes desenvolvam aptidões e atitudes proativas. \textbf{Posto} isso, não se pode esquecer \textbf{que} as formações profissionais a serem oferecidas precisam estar adequadas à demanda local e regional, o \textbf{que} faz com \textbf{que} se reconheçam as diferenças e as respeitem, como sendo balizadoras e propulsoras do crescimento intelectual-dialógico, pois é, por \textbf{intermédio} da escuta proativa, da empatia, do respeito às diversidades, da pluralidade e das interculturas \textbf{que} os \textbf{sujeitos} ativos do processo \textbf{dialógico} efetivam -se como agentes do processo intelectual, social e histórico, demonstrando a capacidade de articular ideias, conceitos, bem como conhecimentos, habilidades e atitudes, isto é, saber, saber fazer e fazer.
Outro fator importante \textbf{que} se deve considerar a partir do Relatório de Escuta ( RS, 2019c ) é a necessidade de investimento em formação e qualificação profissional dos professores, pois ao mesmo tempo \textbf{que} esses profissionais gaúchos não se \textbf{dizem} totalmente preparados para \textbf{amparar} o jovem na construção de \textbf{um} projeto de vida, este jovem espera \textbf{que} seus mestres os auxiliem nesse processo.
Assim, para preparar esse professor, se faz necessário \textbf{um} investimento social ; é \textbf{preciso} \textbf{que} o tempo seja o seu parceiro : tempo para planejar, para mentorear, para adquirir conhecimento e decodificar \textbf{ideais} de vida.
O professor precisa dispor do seu tempo, para \textbf{que} seja instaurada \textbf{uma} formação \textbf{crítico-reflexiva}.
Somente quando esse profissional entender realmente o seu papel na sociedade, a sua relevância na vida do outro, os impactos e condicionantes materiais e simbólicos é \textbf{que}, conforme Bourdieu ( 2007 ), se começará a introduzir o habitus docente, o qual comporta as dimensões material, simbólica e cultural.
É o uso do Corpo-si ( SCHWARTZ ; DURRIVE, 2010 ), é o labor da Experiência ( BONDÍA, 2002 ).
Nesse cenário, os profissionais da educação cumprem \textbf{um} papel indiscutível : a efetivação do Ensino Médio como práxis escolar, como acontecimento no cotidiano da escola.
São esses profissionais \textbf{que}, dotados das suas capacidades de gestão do processo de aprendizagem, dão materialidade ao novo fazer pedagógico indo ao encontro do novo momento da educação no Estado do Rio Grande do Sul.
Dessa forma, o Referencial Curricular Gaúcho para o Ensino Médio deve englobar o atendimento às diferentes demandas das juventudes e dos \textbf{sujeitos} de cada região, dialogando com as diferenças entre as realidades, conectando -as por suas semelhanças e necessidades em comum para, então, consolidar e aprofundar os conhecimentos, garantindo, assim, o aprimoramento humano.

% matched lemmas: abordagem, amparar, aproximar, basear, conhecido, contudo, crítico-reflexivo, demandar, desencadear, dialógico, direção, dizer, enquanto, equidade, exigir, explícito, forte, fundamento, ideal, imprescindível, intermédio, lado, logo, maturidade, novidade, permear, postar, preciso, prova, que, singularidade, sujeito, sumo, sólido, tarefa, um, ver, viver
\end{textsample}
